\section{Síntesis y ampliación}

\subsection{Conclusiones centrales}

\begin{enumerate}
\item La era de los Agent no consiste en una mejora de la capacidad de conversación, sino en que el modelo entra en entornos, herramientas y tareas de largo alcance.
\item El valor de OpenClaw está en que generó un consenso en torno al framework y permitió a los equipos ver que el modelo y el scaffold pueden potenciarse mutuamente.
\item El postentrenamiento, en particular el Agent RL scaling, se convierte en la línea principal, y requiere entornos de tareas, rollout, reward, eval y control de costos.
\item La importancia de MiMo-V2 está en la orquestación de múltiples modelos, no en el lanzamiento de un modelo individual.
\item La asignación de capacidad de cómputo pasa del predominio del preentrenamiento a un reparto más equilibrado entre research, pre-train y post-train.
\item La horizontalidad organizacional, la diversidad y los valores culturales influyen en la rapidez con que un AI Lab adopta un nuevo paradigma.
\item Las oportunidades de los equipos chinos para alcanzar a los líderes están en el Agent post-train, la sensibilidad a los costos y la retroalimentación de las aplicaciones; los retos están en la RL infra y en la agilidad organizacional.
\end{enumerate}

\subsection{Preguntas abiertas}

\begin{itemize}
\item ¿Cuál será la forma predominante de reward y de entorno en el Agent post-train?
\item ¿Seguirá siendo a largo plazo un modelo base de 1T el requisito de entrada, o será desplazado gradualmente por modelos pequeños + framework?
\item ¿La orquestación de múltiples modelos desembocará en un ecosistema unificado o en el enrutamiento de modelos dentro de un mercado abierto?
\item ¿Cómo puede la horizontalidad organizacional evitar convertirse en una colaboración caótica y aumentar de verdad la velocidad de la investigación?
\item Cuando el modelo no tiene una «crisis de supervivencia», ¿cómo deben los seres humanos diseñar su función de valor y su lugar en la sociedad?
\end{itemize}

\subsection{Lecturas adicionales}

\begin{itemize}
\item EP139, Su Yu (苏煜), historia técnica de los Agent: para entender el contexto histórico del OpenClaw Moment y de los Language Agent.
\item Materiales sobre Claude Code / OpenClaw / Computer Use Agent: para entender cómo los frameworks de Agent cambian la forma de los productos.
\item Artículos sobre RLHF, RLAIF, verifier, rollout y agent benchmark: para entender la infraestructura del postentrenamiento.
\item MoE, attention alternatives, multi-modal model orchestration: para entender los problemas de sistemas de modelos detrás de MiMo-V2.
\end{itemize}

\begin{importantbox}{Valoración final}
Quién gane en la era de los Agent no lo decidirá un único modelo aislado. Se parece más a una competencia de sistemas: quien logre integrar más rápido en un ciclo cerrado el modelo, el framework, los entornos de tareas, el postentrenamiento, la capacidad de cómputo, el producto y la organización tendrá más probabilidades de ir a la cabeza en la siguiente etapa.
</br>
\end{importantbox}

\end{document}
